\section{Co-registered and geocoded SLC (CSLC \& GSLC)}
\label{sec:cslc_gslc}

\subsection{Product Description}
\label{sec:cslc_product_description}
Image co-registration refers to the process of aligning two or more images so that each pixel corresponds to the same geographic location in the target scene. This alignment can achieve subpixel accuracy and it is a critical step in generating multidimensional image stacks, enabling advanced data analysis.

For the ROSE-L mission, the co-registered products are provided in radar geometry \cite{RD41}. This product, referred to as the \acs{cslc} (Co-Registered Single-Look Complex), plays a central role in the interferometric process. It incorporates corrections for geometric distortions (e.g., topographic effects), atmospheric delays (tropospheric and ionospheric), and other influencing factors. 

The \acs{cslc} maintains the full fidelity of the original SAR data, introducing no distortions, and serves as a foundational product from which other \acs{ard} products can be derived.
The amplitude of the \acs{cslc} corresponds to \gls{beta0}, obtained by the calibration step using the Look-up-Table from the \acs{slc} annotations.
The phase information in a \acs{cslc} is flattened with respect to the reference orbit and a Digital Elevation Model (DEM).

To support analysis in map-projected (ground) geometry, \acs{slc} or \acs{cslc} products can be automatically resampled onto a common geographic grid (e.g., \acs{utm}, \acs{ups}). This transformation removes the topographic phase and mitigates artifacts related to atmospheric disturbances or orbital inaccuracies. This map-projected product is known as the \acs{gslc} (Geocoded Single-Look Complex). While \acs{gslc} simplifies user interaction, allowing direct visualisation in \acs{gis} tools and the integration of data from multiple satellites, it is inherently less rigorous and unsuitable for high-precision interferometric applications such as coherence analysis or phase unwrapping.
Like the \acs{slc}, the phase of the \acs{gslc} is flattened with respect to the reference orbit and DEM in order to eliminate topographic contributions \cite{Zebker2017, Zheng2017}. The spatial resolution of the \acs{gslc} in map coordinates is comparable to the native resolution of the original \acs{slc}.

The \acs{cslc} product consists of \acs{slc} data that has been resampled to a reference radar grid, which may be:
\begin{itemize}
	\item virtual: based on a nominal reference orbit and sensor parameters,
	\item real: using the radar grid of a specific \acs{slc} image as the reference.
\end{itemize}

Although co-registration is essential for generating most \acs{ard} products and enabling many applications, the \acs{cslc} is often considered an intermediate product in the processing chain. However, it can also serve as a final \acs{ard} product for users who lack the expertise or the resources to perform advanced co-registration and topographic phase removal, but still wish to produce more sophisticated products than a basic \acs{ard}.

The \acs{cslc} is additionally foreseen as a standard \acs{ard} product in the ongoing re-engineering of Sentinel-1 Level-1 and Level-2 processors \cite{S1_RP_portal, S1_RP_ProdSpec}.

This planned product aligns with the evolution of the Sentinel-1 processor chain, ensuring consistency across Copernicus missions as required by the re-engineering project.

The dataset is formatted as a ZARR file. This format supports the inclusion of both the processed measurement data and the associated metadata within the same file, making the product self-contained and self-describing, with no need for additional annotations or supporting files.

\subsubsection{Input and Auxiliary Data}
\label{sec:aux_data_cslc}
The following input and auxiliary data are considered for product generation:

\begin{longtblr}[
	label = tab:in_aux_cslc_products,
	entry = none,
	caption = {\acs{cslc}/\acs{gslc} input and auxiliary data},
	]{
		width = \linewidth,
		colspec = {Q[l,20]Q[l,60]Q[c,20]},
		hlines,
		hline{1-Z} = {-}{0.05em},
	}
	\textbf{Name}   & \textbf{Description}   & \textbf{Mandatory (M) / Optional (O)} \\
	\acs{slc} & L1 Single Look Complex Images and annotations & M \\
    Radar Reference Grid & Virtual or real radar reference grid & M \\
    \acs{dem} & External DEM (e.g., Copernicus DEM at different resolutions, with 30 m as the default worldwide). & M \\
    Orbits & Standard, restituted or precise orbits. & M \\
    \acs{etad} & \acs{etad}-like products for atmospheric delay and timing corrections. & O \\
    \acs{tec} & External \acs{tec} data for delay correction & O \\
\end{longtblr}

\subsubsection{Main Variables}
\label{sec:variables_cslc}
The following tables list the main variables associated with the \acs{cslc} and \acs{gslc} products, respectively.

\begin{longtblr}[
	label = tab:var_cslc_products,
	entry = none,
	caption = {\acs{cslc} main variables},
	]{
		width = \linewidth,
		colspec = {Q[l,15]Q[l,52]Q[l,20]Q[l,13]},
		hlines,
		hline{1-Z} = {-}{0.05em},
	}
	\textbf{Name} & \textbf{Description} & \textbf{Dimensions} & \textbf{Data Type} \\
	cslc & Co-registered single look complex data: 
    
    - the amplitude is in \gls{beta0} convention, obtained by the calibration step using the Look-up-Table from the \acs{slc} annotations;
    
    - the phase corresponds to the interpolated and resampled \acs{slc} flattened phase with respect to the reference orbit and the input Digital Elevation Model (DEM) & polarisation, acquisition\_date, azimuth\_time, range\_time & complex64 \\
    phase\_model & Topographic phase model, in rad, removed from the original image phase & acquisition\_date, azimuth\_time, range\_time & float32 \\
    
\end{longtblr}

\begin{longtblr}[
	label = tab:var_gslc_products,
	entry = none,
	caption = {\acs{gslc} main variables},
	]{
		width = \linewidth,
		colspec = {Q[l,15]Q[l,52]Q[l,20]Q[l,13]},
		hlines,
		hline{1-Z} = {-}{0.05em},
	}
	\textbf{Name} & \textbf{Description} & \textbf{Dimensions} & \textbf{Data Type} \\
	gslc & Projected \acs{slc} or \acs{cslc}, with module equal by default to \gls{beta0} or, optionally to \gls{gamma0}, coefficient and phase equal to the \acs{slc} or \acs{cslc} & polarisation, acquisition\_date, x, y & complex32 \\
    phase\_model & Projected topographic phase model, in rad, removed from the original image phase & acquisition\_date, azimuth\_time, range\_time & float32 \\
    local\_inc
    
    \_angle & DEM-based local Incidence angle image [deg]. This variable is present only if the \acs{gslc} is provided in the \gls{gamma0} convention and is generated from an \acs{slc}. If the \acs{gslc} is derived from a \acs{cslc}, this variable will be null, and the reference\_inc\_angle metadata will be present & x, y & float32 \\
    anf & Area normalisation factor or local contributing area, used to obtain the gamma0 coefficient. This variable is present only if the \acs{gslc} is provided in the \gls{gamma0} convention and is generated from a \acs{slc}. If the \acs{gslc} is derived from a \acs{cslc}, this variable will be null, and the reference\_anf metadata will be present & x, y & float32 \\
\end{longtblr}

Note that, for the \acs{gslc} the fields "anf" and "local\_inc\_angle" are always included: as variables for products directly derived from \acs{slc}, and as metadata for products derived from \acs{cslc}, with links to the corresponding static layer.

The table below describes the dimensions/coordinates associated with the \acs{cslc} and \acs{gslc} main variables.

\begin{longtblr}[
	label = tab:var_cslc_coordinates,
	entry = none,
	caption = {\acs{cslc} and \acs{gslc} main variables coordinates/dimensions},
	]{
		width = \linewidth,
		colspec = {Q[l,20]Q[l,62]Q[l,18]},
		hlines,
		hline{1-Z} = {-}{0.05em},
	}
	\textbf{Name} & \textbf{Description} & \textbf{Data Type} \\
	polarisation & Acquired polarisations & string \\
    acquisition\_date & Image acquisition dates, formatted as YYYY-MM-DD & string \\
    azimuth\_time & Time along satellite track direction & datetime64 \\
    range\_time & Time across satellite track direction & float64 \\
    x & UTM/UPS x (East) coordinates [m] & float64 \\
    y & UTM/UPS y (North) coordinates [m] & float64 \\
\end{longtblr}

\subsubsection{Main Metadata}
\label{sec:metadata_cslc}
The following tables list the main metadata and per-pixel metadata associated with the \acs{cslc} and \acs{gslc} products, respectively.

\begin{longtblr}[
	label = tab:meta_cslc_products,
	entry = none,
	caption = {\acs{cslc} main metadata and per-pixel metadata},
	]{
		width = \linewidth,
		colspec = {Q[l,17]Q[l,38]Q[l,20]Q[l,13]Q[c,12]},
		hlines,
		hline{1-Z} = {-}{0.05em},
	}
	\textbf{Name} & \textbf{Description} & \textbf{Dimensions} & \textbf{Data Type} & \textbf{Per-pixel Metadata}\\
	reference\_slc & Link to input reference SLC & - & string & \\
    reference\_orbit & Link to reference orbit file & - & string & \\
    orbit\_coeffs & Polynomial coefficients describing the orbit associated with each input image belonging to the \acs{cslc} stack & acquisition\_date, axis, degree & float64 & \\
    orbit\_reference\
    
    \_time & Reference time associated with the orbit polynomial coefficients for each input image belonging to the \acs{cslc} stack & acquisition\_date & float64 & \\
    orbit\_type & Type of orbital data used for each of the images belonging to the \acs{cslc} stack: ETAD, POE, POD, RES & acquisition\_date & string & \\
    orbit\_files & Orbital file path list used for each of the images belonging to the \acs{cslc} stack & acquisition\_date & string & \\
    reference\_dem & Link to reference DEM & - & string & \\
    etad\_file & ETAD file used for each image belonging to the CSLC stack. This is optional: it is included only if used; otherwise, it is null & - & string & \\ 
    etad\_azcorr & Extended Timing Annotation Data: azimuth time correction resampled onto the common grid of the reference geometry. This is optional: it is included only if used; otherwise, it is null & acquisition\_date, azimuth\_time, slant\_range\_time & float64 & \ding{51} \\
    etad\_rgcorr & Extended Timing Annotation Data: slant range time correction resampled onto the common grid of the reference geometry. This is optional: it is included only if used; otherwise, it is null & acquisition\_date, azimuth\_time, slant\_range\_time & float64 & \ding{51} \\
    tec & Link to external \acs{tec} data, if used & - & string & \\
    reference\_inc
    
    \_angle & Link to DEM-based local incidence angle map. & - & string & \\
    reference\_anf & Link to reference area normalisation factor or local contributing area. & - & string & \\
    look\_vector & 3-D components radar unit look vector, in an Earth-Centred Earth-Fixed (\acs{ecef}) coordinate system. It consists of unit vectors from antenna to surface pixel. Only for reference geometry. &  (x, y) or (azimuth\_time, range\_time) & float64 & \ding{51} \\
    phase\_demod & Phase to be applied to demodulate each image belonging to the \acs{cslc} stack & acquisition\_date, azimuth\_time, slant\_range\_time & float32 & \ding{51} \\
    sp\_side-band & Side-band used for the application of the split-spectrum technique. It is present only when the split-spectrum method is applied; otherwise, it is set to null. \\ 
    beta0 & LUT of the beta0 coefficients obtained from each input \acs{slc} annotation belonging to the \acs{cslc} stack & acquisition\_date, azimuth\_time, slant\_range\_time & float32 & \\
    gamma0 & LUT of the gamma0 coefficients obtained from each input \acs{slc} annotation belonging to the \acs{cslc} stack & acquisition\_date, azimuth\_time, slant\_range\_time & float32 & \\
    rcs & LUT of the Radar Cross Section (RCS) coefficients obtained from each input \acs{slc} annotation belonging to the \acs{cslc} stack & acquisition\_date, azimuth\_time, slant\_range\_time & float32 & \\
    noise & LUT of thermal noise correction obtained from each input \acs{slc} annotation belonging to the \acs{cslc} stack & acquisition\_date, azimuth\_time, slant\_range\_time & float32 & \\
    crg\_accuracy & Estimated co-registration accuracy in correspondence of each patch of the images belonging to the \acs{cslc} stack & acquisition\_date, azimuth\_time, slant\_range\_time & float32 & \\
    crg\_accuracy
    
    \_bias & Estimated co-registration accuracy bias, computed as the mean value of the co-registration accuracy over all image patches & acquisition\_date & float32 & \\
    crg\_accuracy
    
    \_std & Estimated co-registration accuracy standard deviation, computed as the standard deviation of the co-registration accuracy over all image patches  & acquisition\_date & float32 & \\    
    source\_data
    
    \_attributes & Dictionary containing the attributes associated with each image belonging to the \acs{cslc} stack & - & Dictionary & \\
    x\_sample & Pixel spacing in azimuth direction & - & float32 & \\
    y\_sample & Pixel spacing in slant range direction & - & float32 & \\
\end{longtblr}

\begin{longtblr}[
	label = tab:meta_gslc_products,
	entry = none,
	caption = {\acs{gslc} main metadata and per-pixel metadata},
	]{
		width = \linewidth,
		colspec = {Q[l,17]Q[l,38]Q[l,20]Q[l,13]Q[c,12]},
		hlines,
		hline{1-Z} = {-}{0.05em},
	}
	\textbf{Name} & \textbf{Description} & \textbf{Dimensions} & \textbf{Data Type} & \textbf{Per-pixel Metadata}\\
	cslc/slc & Link to input \acs{cslc} or \acs{slc} & - & string & \\
    reference\_dem & Link to reference DEM & - & string & \\
    etad\_file & Link to ETAD file. This is optional: it is included only if used; otherwise, it is null & - & string & \\ 
    reference\_inc
    
    \_angle & Link to DEM-based local incidence angle map. It is present only if the \acs{gslc} is generated from a \acs{cslc} & - & string & \\
    reference\_anf & Link to reference area normalisation factor or local contributing area. It is present only if the \acs{gslc} is generated from a \acs{cslc} & - & string & \\
    look\_vector & 3-D components radar unit look vector, in an Earth-Centred Earth-Fixed (\acs{ecef}) coordinate system. It consists of unit vectors from antenna to surface pixel. Only for reference geometry. &  (x, y) or (azimuth\_time, range\_time) & float64 & \ding{51} \\
    mask & Mask containing, for each ground cell, the following information:
    
     - 0: no distortions
    
     - 1: layover
    
     - 2: shadow
    
     - 3: undefined & x, y & byte & \ding{51} \\
    nlooks & Number of equivalent looks for the gamma flattening. Present if the \acs{gslc} amplitude is in \gls{gamma0} convention & x, y & float32 & \ding{51} \\
    geometric
    
    \_accuracy
    & Expected geometric accuracy, including eastern and northern geometric accuracy bias and standard deviation. The measures will be performed on a validation dataset and periodically updated. & - & dictionary & \\ 
    east\_geom
    
    \_bias & Expected eastern geometric bias & - & float32 & \\
    east\_geom
    
    \_std & Expected eastern geometric standard deviation & - & float32 & \\
    north\_geom
    
    \_bias & Expected northern geometric bias & - & float32 & \\
    north\_geom
    
    \_std & Measured northern geometric standard deviation & - & float32 & \\
    geometric
    
    \_accuracy
    
    \_radial
    
    \_rmse	& Expected radial RMSE of geometric accuracy & - & float32 & \\
    x\_sample & Pixel spacing in x direction & - & float32 & \\
    y\_sample & Pixel spacing in y direction & - & float32 & \\
\end{longtblr}

The table below describes the dimensions/coordinates associated with the \acs{cslc} and \acs{gslc} main metadata.

\begin{longtblr}[
	label = tab:var_cslc_metadata_coordinates,
	entry = none,
	caption = {\acs{cslc} and \acs{gslc} main metadata coordinates/dimensions},
	]{
		width = \linewidth,
		colspec = {Q[l,20]Q[l,62]Q[l,18]},
		hlines,
		hline{1-Z} = {-}{0.05em},
	}
	\textbf{Name} & \textbf{Description} & \textbf{Data Type} \\
	acquisition\_date & Image acquisition dates, formatted as YYYY-MM-DD & string \\
    azimuth\_time & Time along satellite track direction & datetime64 \\
    range\_time & Time across satellite track direction & float64 \\
    x & UTM/UPS x (East) coordinates [m] & float64 \\
    y & UTM/UPS y (North) coordinates [m] & float64 \\
    axis & x, y, z axis index & int64 \\
    degree & orbit polynomial coefficient degree & int64 \\
\end{longtblr}


\subsubsection{Reference System Tiling and Slicing}
\label{sec:cslc_reference_system_tiling}
\autoref{tab:clsc_gslc_ref_system_slicing} reports the reference system, the slicing, and grid for \acs{cslc} and \acs{gslc} products.


\begin{longtblr}[
	label=tab:clsc_gslc_ref_system_slicing,
	entry = none,
	caption = {Reference system and Slicing for \acs{cslc} and \acs{gslc}},
	]{
		width=\linewidth,
		colspec = {Q[l,15]Q[l,40]Q[l,15]Q[l,30]},
		%vline{2-5} = {-}{},
		%vline{5-5} = {-}{},
		hlines,
		hline{1-Z} = {-}{0.05em},
		rowhead=2,
	}
	\textbf{Product} & \textbf{Reference System} & \textbf{Slicing} & \textbf{Spacing}\\
	\textbf{\acs{cslc}} & Radar: azimuth/slant-range & Burst & Same of the input data \\
	\textbf{\acs{gslc}} & Cartographic: UTM/UPS & Burst & Comparable with radar spacing \\   
	
\end{longtblr}

The ideal solution to ensure interoperability between the ROSE-L and Sentinel-1 Copernicus missions would be to adopt both a common slicing approach and a shared reference radar grid.
While the former is unfeasible—since slicing is determined by burst segmentation—the latter could still be explored.
Given that both missions share the same orbit, it is possible to define a common virtual radar grid, with spacing determined by the higher-resolution sensor. This would allow data from the different sensors to be linked using the coordinates of the reference virtual grid.


\subsection{Product Requirements}
\label{sec:cslc_requirements}

\subsubsection{Accuracy}

A critical requirement for the \acs{cslc} is that it must preserve the intrinsic radiometric and geometric characteristics of the SAR imagery. To verify this, we propose the use of quantitative image quality indicators before and after co-registration.
	
The first two indicators employed are the Peak Sidelobe Ratio (\acs{pslr}) and the Integrated Sidelobe Ratio (\acs{islr}), as previously defined. These metrics are sensitive to distortions introduced by interpolation and geometric transformations. Therefore, by comparing \acs{pslr} and \acs{islr} values computed on point targets or isolated scatterers in both the original and the co-registered SAR images (\acs{cslc} product), we can assess whether the process has introduced significant degradation.
	
Ideally, the \acs{pslr} and \acs{islr} values should remain stable, within an acceptable tolerance, and significant deviations would indicate that the interpolation step or geometric transformation negatively affects image quality and could compromise downstream applications such as InSAR or coherent change detection. 

The accuracy requirements reported in the following tables are intended to be evaluated on a set of isolated scatterers or Corner Reflectors (\acs{cr}s) with a radar cross section (\acs{rcs}) greater than 35~dB.


\begin{longtblr}[
	label = tab:cslc_general_accuracy,
	entry = none,
	caption = {Accuracy Requirements for \acs{cslc} product},
]{
	width = \linewidth,
	colspec = {Q[l,30] Q[l,60] Q[c,20]},
	hlines,
	hline{1-Z} = {-}{0.05em},
	rowhead = 1,
}

\textbf{Metric} & \textbf{Description} & \textbf{Requirement} \\

PSLR degradation (dB) & Standard deviation (1-$\sigma$) of degradation of peak sidelobe level relative to mainlobe peak in both range and azimuth directions & $\leq 0.1$~dB \\

ISLR degradation (dB) & Standard deviation (1-$\sigma$) of degradation of total sidelobe energy relative to mainlobe energy in both range and azimuth directions & $\leq 0.1$~dB \\

Resolution degradation (\%) & Standard deviation (1-$\sigma$) of resolution decreasing expressed as fraction of original resolution in both range and azimuth directions & $\leq 1\%$ \\


\end{longtblr}


Regarding the co-registration errors in both azimuth and range directions, the achievable performance is strongly influenced by both the sensor's operating band (with significant differences between C-band and L-band) and the type of co-registration strategy employed. In particular, three distinct co-registration approaches can be identified:
\begin{itemize}
    \item Model-driven. In this case, accurate orbital data is assumed to be available, but no models are used to correct for atmospheric delays.
    \item Model-driven + atmospheric correction. In addition to precise orbital information, this approach assumes the availability of state-of-the-art atmospheric correction models (e.g., \acs{etad}, JHR GIM), which are used to mitigate the impact of atmospheric propagation effects.
    \item Data-driven. This approach includes fine co-registration techniques based on cross-correlation methods, as well as azimuth correction strategies such as Enhanced Spectral Diversity (\acs{esd}), which exploits the spectral diversity in the overlap areas of consecutive bursts to estimate residual azimuth co-registration errors.
\end{itemize}
To provide a clear overview of the co-registration accuracy requirements under different processing assumptions, the expected performance thresholds are summarized separately for the C-band and L-band sensors. Due to the distinct propagation characteristics and to the sensitivity to atmospheric effects associated with each frequency band, the accuracy requirements differ accordingly. The following tables report the maximum acceptable residual errors in azimuth and range directions for each band, considering various co-registration strategies, including model driven approaches (with and without atmospheric correction) and data-driven methods. 

All accuracy values are expressed as fractions of the nominal resolution in the corresponding dimension, providing a normalized reference for evaluating the effectiveness of the co-registration process. The accuracy thresholds are evaluated at 1-$\sigma$.


\begin{longtblr}[
    label = tab:cslc_coregistration_accuracy_cband,
    entry = none,
    caption = {Coregistration Accuracy Requirements for \acs{cslc} product (C-Band)},
    ]{
        width = \linewidth,
        colspec = {Q[l,30] Q[l,60] Q[c,20]},
        hlines,
        hline{1-Z} = {-}{0.05em},
        rowhead = 1,
    }
    
    \textbf{Metric} & \textbf{Description} & \textbf{Requirement} \\
    
    \SetCell[r=2]{l} Range Accuracy (\% pixel) 
    & Model-Driven (only orbit and \acs{dem}) & $\leq 5\%$\\
    & Model-Driven + atmospheric correction & $\leq 2\%$ \\

    \SetCell[r=2]{l} Azimuth Accuracy (\% pixel) 
    & Model-Driven (only orbit and \acs{dem}) & $\leq 2\%$\\
    & Data-Driven (\acs{esd}) & $\leq 0.5\%$ \\
    
\end{longtblr}

\begin{longtblr}[
    label = tab:cslc_coregistration_accuracy_lband,
    entry = none,
    caption = {Coregistration Accuracy Requirements for \acs{cslc} product (L-Band)},
    ]{
        width = \linewidth,
        colspec = {Q[l,30] Q[l,60] Q[c,20]},
        hlines,
        hline{1-Z} = {-}{0.05em},
        rowhead = 1,
    }
    
    \textbf{Metric} & \textbf{Description} & \textbf{Requirement} \\
    
    \SetCell[r=3]{l} Range Accuracy (\% pixel) 
    & Model-Driven (only orbit and \acs{dem})& $\leq 60\%$\\
    & Model-Driven + atmospheric correction & $\leq 25\%$ \\
    & Model-Driven + atmospheric correction + Data-Driven (\acs{icc}) & $\leq 10\%$ \\

    \SetCell[r=3]{l} Azimuth Accuracy (\% pixel) 
    & Model-Driven (only orbit and \acs{dem}) & $\leq 150\%$ \\
    & Model-Driven + atmospheric correction & $\leq 50\%$ \\
    & Data-Driven (\acs{icc}) & $\leq 10\%$ \\
    
\end{longtblr}
In the case of L-band SAR systems, atmospheric effects -in particular those related to ionospheric disturbances- can significantly impact the accuracy of fine co-registration. As a result, the achievable co-registration performance is expected to exhibit strong variability depending on several factors, including the geographic location of the acquisition, the time of day, and prevailing space weather conditions.

The co-registration accuracy requirements reported herein are defined with respect to worst-case scenarios, such as acquisitions conducted at high latitudes (e.g., near the polar regions) or during periods of elevated solar activity, where ionospheric perturbations are more pronounced. Nevertheless, in more favourable conditions —such as acquisitions in equatorial or mid-latitude regions during geomagnetically quiet periods— substantially improved performance is anticipated, potentially exceeding the stated requirements by up to an order of magnitude.

\begin{comment}
For main variables and per-pixel metadata (if applicable): Schematic table with accuracy depending on the possible input datasets
\begin{itemize}
    \item IRF degradation wrt SLC (tool SCT) to be evaluated on few S1 \& SAOCOM data
    \item PSLR degradation wrt SLC (tool SCT) to be evaluated on few S1 \& SAOCOM data 
    \item ... degradation wrt SLC (tool SCT) to be evaluated on few S1 \& SAOCOM data
    \item geolocation accuracy (tool SCT) in correspondence of CRS
    \item Spectral information preservation (doppler difference wrt input data)
    \item Coregistration accuracy (rg, az) (theoretical formulas, CRS?) as a function of Ancillary data / processing
    \item phase continuity between bursts (theoretical analysis) ???? as a function of Ancillary data / processing
\end{itemize}
\end{comment}

\subsubsection{Timeliness}
\label{sec:cslc_timeliness}

In principle \acs{cslc} product could be generated as soon as L1 data are available. The quality and the accuracy of the co-registration process strongly depends on the availability of precise orbital data, models of the atmosphere (troposphere and ionosphere). One could foresee multiple releases of the product, with successive corrections according to the availability of precise orbits and external models. A first release would be within a few hours after the acquisition, for time-critical applications. A first revision including precise orbital data could be provided after few days, and the final release could be delivered whenever numerical weather reanalysis data, ETAD and accurate ionospheric models are available.

 In this approach the only criticality is on the availability of precise orbital data (in general available in about one week). \autoref{fig:cslc_basic} shows the flowchart for the generation of \acs{cslc} product.

In addition, the provided timeline indicates the time required to collect the input data and complete the main processing steps. The processing time refers to a standard \acs{s1} dataset and is approximated, considering that the latency in obtaining the input data is the dominant factor.

\begin{figure}[H]
	\centering
	\begin{tikzpicture}[node distance=1cm and 2cm]
		
		% Define the style for the nodes
		% \tikzstyle{node} = [circle, draw, fill=blue!20, text centered, minimum size=1.5cm]
		
		% Draw nodes
		\node[rectin] (init) {\acs{slc}};
		\node[block, right=of init] (coreg_geo) {Geometric co-registration};
		\node[rectaux, below=of init] (dem) {\acs{dem}};
            \node[empty, below=of coreg_geo] (empty1) {};
		\node[rectaux, below=of empty1] (orbit) {Orbital data};
		\node[rectout, right=of coreg_geo] (cslc) {\acs{cslc}};
		
		% Draw edges with arrows
		\draw[->, thick] (init) -- (coreg_geo);
		\draw[->, thick] (dem) -| (coreg_geo);
		\draw[->, thick] (orbit) -- (coreg_geo);
		\draw[->, thick] (coreg_geo) -- (cslc);
        % Timeline with linked markers and adaptive length
        \def\timelineStart{init}
        \def\timelineEnd{cslc}
        \draw[thick, ->] (\timelineStart.west |- 0,-6) -- (\timelineEnd.east |- 0,-6) node[right] {Time};
        \foreach \node/\label in {init/Start, orbit/1h-3d, cslc/<1h} {
            \draw (\node.south |- 0,-6) -- (\node.south |- 0,-6.5) node[below] {\label};
            \draw[dotted, thick] (\node.south) -- (\node.south |- 0, -6.25); % Connects the node to the timeline
        }
    
	\end{tikzpicture}
	\caption{Sample workflow for generation of \acs{cslc} using geometric co-registration method.}
    \label{fig:cslc_basic}
\end{figure}

To further improve product quality, additional data, as listed in \autoref{sec:aux_data_cslc}, can be used. This approach introduces a delay in data production, which depends on the availability of external models: \acs{etad}-like, troposphere and ionosphere. \autoref{fig:cslc_ref_mdl} shows the flowchart for refinement of co-registration considering external models.

\begin{figure}[H]
	\centering
	\resizebox{\textwidth}{!}{
	\begin{tikzpicture}[node distance=1cm and 2cm]
		
		% Define the style for the nodes
		% \tikzstyle{node} = [circle, draw, fill=blue!20, text centered, minimum size=1.5cm]
		
		% Draw nodes
		\node[rectin] (init) {\acs{slc}};
		\node[block, right=of init] (coreg_geo) {Geometric Co-registration};
		\node[block, right=of coreg_geo] (coreg_ref) {Refinement};
		\node[rectaux, below=of init] (dem) {\acs{dem}};
		\node[empty, below=of coreg_geo] (empty1) {};
            \node[rectaux, below=of empty1] (orbit) {Orbital data};
		\node[empty, below=of coreg_ref] (empty2) {};
		\node[rectaux, below=of empty2] (models) {External models};
		\node[rectout, right=of coreg_ref] (cslc) {\acs{cslc}};
		
		% Draw edges with arrows
		\draw[->, thick] (init) -- (coreg_geo);
		\draw[->, thick] (coreg_geo) -- (coreg_ref);
		\draw[->, thick] (dem) -| (coreg_geo);
		\draw[->, thick] (orbit) -- (coreg_geo);
		\draw[->, thick] (models) -- (coreg_ref);
		\draw[->, thick] (coreg_ref) -- (cslc);
            % Timeline with linked markers and adaptive length
            \def\timelineStart{init}
            \def\timelineEnd{cslc}
            \draw[thick, ->] (\timelineStart.west |- 0,-6) -- (\timelineEnd.east |- 0,-6) node[right] {Time};
            \foreach \node/\label in {init/Start, orbit/1h-3d, models/\textbf{1-30d}, cslc/<1h} {
                \draw (\node.south |- 0,-6) -- (\node.south |- 0,-6.5) node[below] {\label};
                \draw[dotted, thick] (\node.south) -- (\node.south |- 0, -6.25); % Connects the node to the timeline
            }
    
	\end{tikzpicture}
	}
	\caption{Sample workflow of generation of \acs{cslc} using geometric co-registration method plus co-registration refinement with external models.}
    \label{fig:cslc_ref_mdl}
\end{figure}


Alternatively, a data-driven approach can be followed, estimating corrections directly from a stack of SAR data. In this case, the processing is more complex, generally involving more than a single image and requiring the solution of inverse problems with high computational complexity.

\autoref{fig:cslc_data_driven} shows the flowchart for the generation of \acs{cslc} product, considering data-driven co-registration method.

\begin{figure}[H]
	\centering
	\resizebox{\textwidth}{!}{
	\begin{tikzpicture}[node distance=1cm and 2cm]
		
		% Define the style for the nodes
		% \tikzstyle{node} = [circle, draw, fill=blue!20, text centered, minimum size=1.5cm]
		
		% Draw nodes
		\node[rectin] (init) {\acs{slc}};
		\node[block, right=of init] (coreg_geo) {Geometric co-registration};
		\node[block, right=of coreg_geo] (coreg_data_driven) {Data-driven Co-registration};
		\node[rectaux, below=of init] (dem) {\acs{dem}};
		\node[empty, below=of coreg_geo] (empty1) {};
            \node[rectaux, below=of empty1] (orbit) {Orbital data};
		\node[empty, below=of coreg_data_driven] (empty2) {};
		\node[rectin, below=of empty2] (ref_stack) {Reference SLC stack};
		\node[rectout, right=of coreg_data_driven] (cslc) {\acs{cslc}};
		
		% Draw edges with arrows
		\draw[->, thick] (init) -- (coreg_geo);
		\draw[->, thick] (coreg_geo) -- (coreg_data_driven);
		\draw[->, thick] (dem) -| (coreg_geo);
		\draw[->, thick] (orbit) -- (coreg_geo);
		\draw[->, thick] (ref_stack) -- (coreg_data_driven);
		\draw[->, thick] (coreg_data_driven) -- (cslc);
            % Timeline with linked markers and adaptive length
            \def\timelineStart{init}
            \def\timelineEnd{cslc}
            \draw[thick, ->] (\timelineStart.west |- 0,-6) -- (\timelineEnd.east |- 0,-6) node[right] {Time};
            \foreach \node/\label in {init/Start, orbit/1h-3d, ref_stack/<1h, cslc/<1h} {
                \draw (\node.south |- 0,-6) -- (\node.south |- 0,-6.5) node[below] {\label};
                \draw[dotted, thick] (\node.south) -- (\node.south |- 0, -6.25); % Connects the node to the timeline
            }
	\end{tikzpicture}
	}
	\caption{Sample workflow for generation of \acs{cslc} using geometric co-registration method plus co-registration refinement exploiting data-driven approach considering a \acs{cslc} stack.}
    \label{fig:cslc_data_driven}
\end{figure}


\subsection{Usability, Integrability and Interoperability}
\label{sec:usability_integrability_CSLC}
The characteristics of the \acs{cslc} and \acs{gslc} products in terms of usability, integrability and interoperability are in the following tables.

\begin{longtblr}[
  caption = {Usability, Integrability, and Interoperability of \acs{cslc} product},
  label = {tab:cslc_usability_integrability_interoperability},
]{
  colspec = {Q[2cm] Q[12cm]},
  rowhead = 1,
  width = \linewidth
}
\hline
\textbf{Aspect} & \textbf{Description} \\
\hline
Usability &
\begin{itemize}
  \item Naturally suited for time-series analysis (e.g., surface deformation, change detection, etc.)
  \item Preserves SAR geometry for reliable spectral and filtering operations
  \item Enables phase-based and amplitude-based applications like InSAR and pixel offset tracking
\end{itemize}
\\
Integrability &
\begin{itemize}
  \item Acts as a core intermediate product in SAR processing chains (e.g., interferometry, change detection)
  \item Easily combined in multi-temporal and joint analyses
\end{itemize}
\\
Interoperability &
\begin{itemize}
  \item Cross-sensor compatibility for ROSE-L and Sentinel-1 missions
  \item Works with common \acs{eo} software
\end{itemize}
\\
\hline
\end{longtblr}

\begin{longtblr}[
  caption = {Usability, Integrability, and Interoperability of \acs{gslc} product},
  label = {tab:gsclc_characteristics},
]{
  colspec = {Q[2cm] Q[12cm]},
  rowhead = 1,
  width = \linewidth
}
\hline
\textbf{Aspect} & \textbf{Description} \\
\hline
Usability &
\begin{itemize}
  \item Geocoded and projected to ground range \acs{cslc} or \acs{slc}, making it more intuitive and accessible for non-expert users
  \item Enables direct visual interpretation and integration with other map-based or \acs{gis} data
  \item When geocoded with adequate accuracy, phase data is suitable for generating interferograms directly on the ground
  \item If radiometrically corrected, the \acs{gslc} amplitude can carry the same information as the \acs{nrb} product
\end{itemize}
\\
Integrability &
\begin{itemize}
  \item Directly usable in \acs{gis} platforms, map viewers, or decision-support tools
  \item Seamlessly integrates into workflows for change detection, urban monitoring, land deformation, and disaster response
  \item If radiometric correction is applied, it can substitute or complement the \acs{nrb} in joint or thematic analyses
\end{itemize}
\\
Interoperability &
\begin{itemize}
  \item High interoperability due to standard geographic referencing and pixel alignment across scenes
  \item Compatible with multi-source datasets (e.g., optical imagery, elevation models, land cover maps)
  \item Adheres to ARD principles, enhancing integration with downstream services and platforms
  \item Facilitates cross-mission usage through harmonized ground projection
  \item Compatible with common \acs{eo} frameworks and pipelines using georeferenced data
\end{itemize}
\\
\hline
\end{longtblr}

